\documentclass[11pt,a4paper]{article}
\input{tutorial_style}

\begin{document}
\tutorialtitle{Note 1 Implementation Companion}{Python Examples for Cyber Control, Differential Games, RL, and MADRL}{From mathematical formulation to reproducible Note 1 experiments}{Companion to Note 1: Cyber Control and Game Learning}
\tableofcontents
\newpage

\section{Purpose}
This companion explains the Note 1 Python files. It maps equations to runnable code and points to extension hooks used in research prototypes. PINN/PIDL examples are handled in the separate Note 2 repository.

\section{File Map}
\begin{longtable}{p{0.28\textwidth}p{0.62\textwidth}}
\toprule
File & Role \\
\midrule
\small\texttt{shared\_setup.py} & local helper that finds the shared foundation package in a sibling workspace \\
\small\texttt{cyber\_dynamics.py} & compatibility wrapper for shared \texttt{cybercontrol} ODE and RK4 utilities \\
\small\texttt{sampled\_continuous}\newline\small\texttt{\_impulse\_env.py} & sampled-data environment with continuous flow, jump map, and parameterized actions \\
\small\texttt{fbsm\_malware\_baseline.py} & open-loop FBSM baseline from PMP \\
\small\texttt{ddqn\_cyber\_defense.py} & discrete-action DDQN defender against a scripted attacker \\
\small\texttt{madrl\_ctde}\newline\small\texttt{\_parameterized\_game.py} & compact CTDE attacker-defender baseline \\
\small\texttt{node\_level\_robustness.py} & node-level epidemic-model robustness experiment \\
\bottomrule
\end{longtable}

\section{General Implementation Architecture}
A clean implementation should separate four layers:
\begin{enumerate}
\item model layer: ODE right-hand side, parameters, constraints;
\item environment layer: reset, action decoding, jump map, ODE integration, reward;
\item algorithm layer: FBSM, DDQN, and MADRL;
\item evaluation layer: baselines, robustness, plots, logs, seeds.
\end{enumerate}
The code folder follows this structure. The environment is intentionally independent of DQN, PPO, or MADRL so that different learning algorithms can reuse the same cyber model.

\section{ODE, Sampled Flow, and Impulse Environment}
The central environment step is
\begin{equation}
\boxed{s_k}\rightarrow\boxed{a_D^k,a_A^k}\rightarrow\boxed{x_k^+=G(x_k,a_k)}\rightarrow\boxed{x_{k+1}=\operatorname{RK4}(f,x_k^+)}\rightarrow\boxed{r_D^k,r_A^k}.
\end{equation}
The implementation is in \texttt{sampled\_continuous\_impulse\_env.py}. The most important design choice is to keep jump effects separate from continuous flow effects.

\begin{lstlisting}[style=pythonstyle,caption={Core idea of a sampled-data flow and impulse environment step.}]
def step(self, defender_action, attacker_action):
    dpar = self.defense_parameters(defender_action)
    apar = self.attack_parameters(attacker_action)
    pre_jump = self.state.copy()
    post_jump = self.jump_map(pre_jump, dpar)
    rhs = lambda x, t: sampled_sir_flow_rhs(x, dpar, apar, self.cfg.params)
    next_state, path = rk4_integrate(rhs, post_jump, t0=self.t*self.cfg.dt,
                                     dt=self.cfg.dt, substeps=self.cfg.substeps,
                                     project=project_simplex3)
    reward = compute_reward(path, dpar, apar)
    return next_state, reward, done, info
\end{lstlisting}

\subsection{How to extend this environment}
To extend from the simple $[S,I,R]$ model to a degree-based network model, replace the state vector with
\begin{equation}
x=[S_1,I_1,R_1,\ldots,S_K,I_K,R_K]
\end{equation}
where $K$ is the number of degree classes. Then update the RHS so that infection pressure for degree class $k$ uses a degree-weighted infected neighbor probability. Deception remains an action-dependent reduction of effective infection pressure during a decision interval rather than a core state compartment. The RL environment interface does not change.

\section{FBSM Baseline}
The FBSM file solves a continuous-control malware problem. It should be used as a classical baseline when the model is known and the control is open-loop. It is not a substitute for DDQN or MADRL because it does not learn a feedback policy.

\begin{lstlisting}[style=pythonstyle,caption={FBSM update logic.}]
for iteration in range(max_iter):
    # forward state integration using current u(t)
    for k in range(n):
        x[k+1] = rk4_state_step(x[k], u[k], h, beta, gamma)
    # backward costate integration from terminal condition
    lambda[n] = terminal_gradient
    for k in range(n, 0, -1):
        lambda[k-1] = rk4_costate_step_backward(x[k], lambda[k], u[k], h)
    # stationarity update and relaxation
    u_new[k] = clip(S[k]*(lambda_S[k]-lambda_R[k])/B, 0, u_max)
    u = relax*u_new + (1-relax)*u
\end{lstlisting}

\section{DDQN Example}
The DDQN example trains the defender with discrete actions. The attacker is scripted so that the defender first learns against a stable environment. This is the recommended first step before full MADRL.

\begin{workflowbox}
\textbf{DDQN workflow.}
\begin{enumerate}
\item initialize Q-network and target Q-network;
\item collect transitions $(s,a,r,s',d)$ using epsilon-greedy exploration;
\item store transitions in replay buffer;
\item sample mini-batches;
\item compute DDQN target;
\item update the online network;
\item periodically copy online weights to target weights;
\item validate against static and rule-based policies.
\end{enumerate}
\end{workflowbox}

\section{MADRL / CTDE Example}
The MADRL file uses decentralized actors and centralized critics. It is intentionally compact. The defender and attacker both receive the same state observation in the example, but this can be changed to partial observations.

\subsection{General extension to richer MADRL}
For research runs, add:
\begin{itemize}
\item generalized advantage estimation (GAE);
\item clipped PPO ratios for MAPPO;
\item opponent pools to reduce cycling;
\item multiple seeds and cross-play matrices;
\item best-response retraining as an approximate exploitability test;
\item separate observation functions for attacker and defender.
\end{itemize}

\section{Evaluation Scripts to Add in Future Work}
The current scripts print losses and returns. A research version should save:
\begin{itemize}
\item CSV trajectories for $S,I,R$ together with action, reward, and baseline diagnostics;
\item training curves;
\item baseline comparison tables;
\item robustness results across parameter perturbations;
\item policy action distributions over time.
\end{itemize}

\section{Complex Extensions}
\subsection{Degree-based HODGE-style model}
To implement a degree-based propaganda model, let the state contain $P_k$ and $C_k$ for each degree $k$. The force of influence uses
\begin{equation}
\Theta_P(t)=\frac{\sum_k k p_k P_k(t)}{\sum_k k p_k},\qquad
\Theta_C(t)=\frac{\sum_k k p_k C_k(t)}{\sum_k k p_k}.
\end{equation}
Continuous propaganda investment changes ODE rates, while impulsive counter-propaganda applies a jump map at selected epochs. The same environment structure still applies.

\subsection{Connection to Note 2}
If the research question becomes inverse learning, missing-mechanism learning, direct neural control, or PMP-informed PINNs, continue in the Note 2 repository. That repository contains the PINN/PIDL code, loss terms, and neural-network checks. Note 1 should remain focused on continuous-control baselines, sampled-data MDPs, and attacker-defender Markov games.

\section{Suggested Project Sequence}
\begin{enumerate}
\item Run FBSM and DDQN on the same simple malware model.
\item Add parameterized actions and compare no defense, static patch, rule-based, DDQN.
\item Add attacker learning and compare independent learning with CTDE.
\item Add node-level graph experiments and parameter-mismatch robustness.
\item Move to Note 2 if the next step is inverse PINN, PIDL, or PMP-informed neural control.
\item Extend to APT or misinformation dynamics after the simpler timing and baseline checks are stable.
\end{enumerate}

\end{document}
